\subsection{Multi-Task Generalization}
To demonstrate the robustness of DeNovo across diverse ADMET endpoints, we evaluated our Attention-GIN encoder on the ClinTox (toxicity) and BBBP (Blood-Brain Barrier Penetration) datasets. Table \ref{tab:multitask} summarizes these results.

\begin{table}[h]
\centering
\caption{Performance Across ADMET Tasks}
\label{tab:multitask}
\begin{tabular}{lcccc}
\toprule
\textbf{Dataset} & \textbf{Task Type} & \textbf{ROC-AUC} & \textbf{Accuracy} & \textbf{Sensitivity} \\
\midrule
Tox21 & Toxicity & 0.837 & 0.812 & 0.765 \\
BBBP & Distribution & 0.825 & 0.892 & 0.915 \\
ClinTox & Toxicity & 0.816 & 0.942 & 0.884 \\
\bottomrule
\end{tabular}
\end{table}

The model achieves consistently high performance ($>0.8$ ROC-AUC) across all tasks. Notably, the high sensitivity on BBBP (0.915) indicates the system is particularly effective at identifying permeable compounds, a critical feature for CNS drug discovery. On ClinTox, the high accuracy (0.942) reflects the model's ability to distinguish safe drugs from those that failed clinical trials, validating its utility in late-stage safety assessment.

\subsection{Faithfulness \& Safety Quantified}
Our faithfulness protocol rejected 36.5\% of LLM-generated explanations for Tox21. In contrast, an unconstrained ablation of our system (removing the verification loop) accepts 100\% of explanations. Given that 36.5\% of these were proven causally inconsistent by our validator, the unconstrained baseline effectively delivers a hallucination rate of over one-third, presenting a severe safety risk. DeNovo's "reject-and-retry" mechanism mitigates this, ensuring that the final output presented to the chemist is causally anchored to the model's prediction.
